Cierta masa gaseosa sufre una transformación al absorber una cantidad de calor $Q$, realizar un trabajo $T$ y sufrir una variación $\Delta U$ en su energía interna. Señale, entre las afirmaciones siguientes, las que son correctas:
\begin{enumerate}
    \item $T = Q$ si la transformación fuera isotérmica.
    \item $\Delta U = Q$ si la transformación fuera isométrica.
    \item $\Delta U = 0$ si la transformación fuera adiabática.
    \item $Q > T$ si la transformación fuera una expansión isobárica.
    \item $Q=0$ si la transformación fuera isotérmica.
\end{enumerate}